\notechapter{N-20250301}{Projective Schemes and Gluing: A Gentle View of Proj}

\section{Why projective schemes are needed}

Projective varieties are built from homogeneous equations. Schemes ask for the same idea at the level of rings, prime ideals, and local functions.

The construction is called \(\operatorname{Proj}\). It is the projective companion of \(\operatorname{Spec}\).

If \(S\) is a graded ring
\[
  S=S_0\oplus S_1\oplus S_2\oplus\cdots,
\]
then \(\operatorname{Proj}S\) is built from homogeneous prime ideals not containing the irrelevant ideal.

This sounds technical, but the basic idea is familiar: projective geometry ignores common scaling, so only homogeneous data make sense.

\section{Projective space}

The most important example is
\[
  \mathbb P^n_R=\operatorname{Proj}R[x_0,\ldots,x_n],
\]
where the variables have degree \(1\).

The standard affine chart \(x_i\ne0\) is written
\[
  D_+(x_i).
\]
On this chart, ratios
\[
  \frac{x_j}{x_i}
\]
become ordinary affine coordinates.

So projective space is assembled from affine pieces, one for each coordinate that is nonzero.

\section{Homogeneous ideals}

A projective subscheme of \(\mathbb P^n\) is described by a homogeneous ideal
\[
  I\subset R[x_0,\ldots,x_n].
\]
Homogeneous equations are necessary because the point
\[
  [x_0:\cdots:x_n]
\]
is unchanged by scaling all coordinates.

For example, a projective conic is given by a homogeneous quadratic equation
\[
  F(x,y,z)=0.
\]
In the chart \(z\ne0\), it becomes an affine equation in \(x/z\) and \(y/z\).

\section{Gluing affine charts}

The most concrete way to think about \(\operatorname{Proj}\) is gluing. Each \(D_+(x_i)\) is affine. The overlaps \(D_+(x_i)\cap D_+(x_j)\) are also affine, and the coordinate changes are given by ratios.

For \(\mathbb P^1\), there are two charts:
\[
  U_0=\{x_0\ne0\},\qquad U_1=\{x_1\ne0\}.
\]
On \(U_0\), use coordinate \(t=x_1/x_0\). On \(U_1\), use coordinate \(s=x_0/x_1\). On the overlap,
\[
  s=\frac1t.
\]
This is the same gluing pattern as the Riemann sphere.

\section{Projective closure revisited}

Given an affine equation \(f(x,y)=0\), homogenizing it gives a projective equation
\[
  F(x,y,z)=0.
\]
The projective scheme remembers not only the affine curve but also what happens at infinity and, if present, nilpotent or embedded structure.

This is a scheme-level version of a familiar geometric act: complete the affine picture by adding the missing directions at infinity.

\section{Why this is not too abstract}

The point of \(\operatorname{Proj}\) is not to make projective geometry harder. It lets projective geometry use the same local-function language as \(\operatorname{Spec}\).

The repeated pattern is:
\[
\begin{array}{c}
  \text{build local affine pieces,}\\
  \text{understand their rings of functions,}\\
  \text{glue them along overlaps.}
\end{array}
\]

\section{Relation with the earlier conic note}

The earlier Erlangen-style conic note emphasized transformations and invariants. The scheme viewpoint emphasizes functions and local rings. These are not competing views.

\[
\begin{array}{c|c}
\text{projective geometry} & \text{incidence, duality, transformations}\\
\text{algebraic geometry} & \text{equations, coordinate rings, local functions}\\
\text{scheme theory} & \text{nilpotents, gluing, base rings}
\end{array}
\]

The same shape becomes richer as the language becomes richer.

%% expanded-on-review

\section{The two-chart construction of the projective line}

The most useful model of \(\operatorname{Proj}\) is still the elementary construction of \(\mathbb P^1\). Take two affine lines:
\[
  U_0=\operatorname{Spec}R[t],\qquad U_1=\operatorname{Spec}R[s].
\]
On the open subsets where \(t\ne0\) and \(s\ne0\), glue by
\[
  s=\frac1t.
\]
The result is projective line.

This example should be kept in mind whenever \(\operatorname{Proj}\) feels abstract. The construction is not a single mysterious space; it is affine geometry plus gluing rules.

\section{Standard opens in the projective plane}

Projective plane has three standard affine charts:
\[
  D_+(x),\qquad D_+(y),\qquad D_+(z).
\]
On \(D_+(z)\), the affine coordinates are
\[
  X=\frac{x}{z},\qquad Y=\frac{y}{z}.
\]
On \(D_+(x)\), one uses
\[
  U=\frac{y}{x},\qquad V=\frac{z}{x}.
\]
The projective plane is what appears after gluing these charts along their overlaps.

This makes projective geometry feel less like adding a mystical line at infinity and more like coordinating several affine windows.

\section{A conic across charts}

The projective conic
\[
  xz-y^2=0
\]
looks different in different charts. On \(z\ne0\), it becomes
\[
  X=Y^2.
\]
On \(x\ne0\), it becomes
\[
  V=U^2.
\]
The same projective curve has multiple affine faces. A single affine picture can hide part of the curve, especially at infinity.

\section{What Proj adds to the earlier story}

The earlier projective note explained homogeneous coordinates and incidence. This note adds the function-theoretic layer:
\[
\begin{array}{c|c}
\text{homogeneous coordinates} & \text{homogeneous primes}\\
\text{affine charts} & \text{standard affine opens}\\
\text{projective equations} & \text{homogeneous ideals}
\end{array}
\]
The new language is heavier, but the basic picture is the same: projective spaces are made by gluing affine charts.


\section{What to remember}

\(\operatorname{Spec}\) turns rings into affine geometric spaces. \(\operatorname{Proj}\) turns graded rings into projective geometric spaces.

\[
  \text{Proj is projective geometry rebuilt from homogeneous algebra and affine gluing.}
\]
This is a good endpoint for the Geometry volume: it returns to projective geometry, but now with the function-theoretic language of schemes.
